%-----------------------------------------------------------------------
\section{Modeling}
%-----------------------------------------------------------------------

Table~\ref{tab:components} provides a consolidated summary of key metrics across
the perception stages. Rather than presenting exhaustive implementation details,
we highlight the aspects that most affect downstream analysis.

\begin{table*}[t]
  \caption{Per-stage performance metrics across the \method{} pipeline.}
  \label{tab:components}
  \centering
  \small
  \begin{tabular}{llr}
    \toprule
    Component & Metric & Result \\
    \midrule
    SA-CNN rally filtering & Validation F1 & 0.903 \\
    SA-CNN & Annotated-hit coverage & 97.32\% \\
    TrackNetV3 tracking & Visible-frame coverage (denoised) & 90.99\% \\
    BST classifier & Validation macro-F1 & 0.805 \\
    YOLO-pose features & Top / Bottom pose-valid rate & 54.29\% / 98.96\% \\
    \bottomrule
  \end{tabular}
\end{table*}

\paragraph{Rally Scene Filtering.}
SA-CNN is a binary classifier that separates court-view frames from non-play
footage—score screens, replays, crowd shots. Trained on 20k labelled images, it
achieves an F1 score of 0.903 and retains 97.32\% of annotated hits, ensuring
that shuttle tracking is not run on irrelevant frames.

\paragraph{Shuttle Tracking.}
TrackNetV3 predicts frame-level shuttle coordinates. Post-processing removes
implausible jumps and interpolates short gaps, achieving a denoised visible
coverage of 90.99\%. Representative trajectories are coherent across rally
frames, although some noise near rally boundaries underscores the need for
precise filtering.

\paragraph{Pose Estimation.}
YOLO-pose extracts body keypoints and positions for both players in each window.
Valid keypoints are obtained for 98.96\% of near-side frames, but only 54.29\%
for far-side frames. This asymmetry proves to be the main bottleneck: when the
far-side pose is absent, classification accuracy drops sharply.

\paragraph{Stroke Classification.}
The Badminton-Stroke Transformer (\texttt{BST\_CG\_AP}) is fine-tuned on
ShuttleSet training clips. On the clean unseen set, it achieves 69.04\%
side-aware top-1 accuracy and 66.64\% macro-F1 across 3{,}178 evaluated strokes.
Player-side accuracy reaches 93.05\%, indicating that the model reliably
identifies which side struck the shuttle even when stroke-type distinctions are
confused.  The pipeline’s principal improvement over earlier versions is the
dramatic reduction of \texttt{none} predictions, from 56.43\% to 6.70\%, which
restores sequential coherence and is critical for the mining analyses.

Table~\ref{tab:per_video} breaks down results across the four held-out
matches. Accuracy ranges from 53.3\% (Video~43, Marin vs.\ Chochuwong
SF) to 76.6\% (Video~44, Antonsen vs.\ Axelsen Finals), a 23-point
spread that reflects match-level variation in broadcast angle and
far-side player visibility. The Finals match (Video~44) benefits from
more consistent near-court exchanges, whereas the semi-final
(Video~43) involves more aggressive backcourt play, compounding
far-side pose errors. Video~41 serves as the reference split used for
validation during training; it is included here only for completeness.

\begin{table}[t]
  \caption{Per-video stroke-classification results on the held-out
  unseen set (videos~41--44). Top-3 accuracy counts a prediction
  correct if the true label appears among the model’s three highest
  scoring classes.}
  \label{tab:per_video}
  \centering
  \small
  \begin{tabular}{clrrrr}
    \toprule
    Video & Match (BWF WTF 2020) & $n$ & Top-1 & Top-3 & Macro-F1 \\
    \midrule
    41 & Ng Ka Long vs.\ Srikanth (QF)  & 1{,}151 & 73.7\% & 90.4\% & 0.723 \\
    42 & Sindhu vs.\ Chochuwong (QF)    &   619   & 64.0\% & 85.1\% & 0.612 \\
    43 & Marin vs.\ Chochuwong (SF)     &   550   & 53.3\% & 74.2\% & 0.478 \\
    44 & Antonsen vs.\ Axelsen (Finals) &   858   & 76.6\% & 95.2\% & 0.725 \\
    \midrule
    \textbf{All} & & \textbf{3{,}178} & \textbf{69.0\%} & \textbf{87.9\%} & \textbf{0.664} \\
    \bottomrule
  \end{tabular}
\end{table}

Table~\ref{tab:top_confusions} lists the most frequent stroke-type
confusion pairs on the test set (side-agnostic). The dominant errors
cluster around \emph{net-zone strokes}: net shot is confused with return
net (89 times) and predicted as \texttt{none} (84 times), and lob is
confused with push (93 times). These three pairs alone account for
266~of the 984 total errors. Crucially, net~shot and return~net are
tactically adjacent---both are soft forecourt exchanges---and are
visually similar when the far-side player is absent or blurred.
The smash$\leftrightarrow$drop confusion (55+49) reflects the shared
overhead mechanics and limited shuttle visibility at the contact point.

\begin{table}[t]
  \caption{Top stroke-type confusion pairs on the 3{,}178-event unseen
  set (side-agnostic; \texttt{none} pairs included where notable).
  Each row is one confusion direction; the symmetric counterpart is
  listed separately when it is also a top error.}
  \label{tab:top_confusions}
  \centering
  \small
  \begin{tabular}{llr}
    \toprule
    True stroke & Predicted as & Errors \\
    \midrule
    lob         & push              & 93 \\
    net shot    & return net        & 89 \\
    net shot    & \texttt{none}     & 84 \\
    drop        & smash             & 55 \\
    smash       & drop              & 49 \\
    push        & lob               & 44 \\
    return net  & drive             & 36 \\
    lob         & clear             & 35 \\
    \bottomrule
  \end{tabular}
\end{table}
